% !PRINTED 40
\weberpage{40}
%<<CONT

\[
(x-\alpha_1)(x-\alpha_2)=x^2-(\alpha_1+\alpha_2)x+\alpha_1\alpha_2,
\]
\[
\begin{aligned}
 (x-\alpha_1)(x-\alpha_2)(x-\alpha_3)&=x^3-(\alpha_1+\alpha_2+\alpha_3)x^2\\
&\quad +(\alpha_1\alpha_2+\alpha_1\alpha_3+\alpha_2\alpha_3)x-\alpha_1\alpha_2\alpha_3.
\end{aligned}
\]
If we assume the law of formation proved for $n-1$ factors, and put
\[
 (x-\alpha_1)(x-\alpha_2)\cdots(x-\alpha_{n-1})=x^{n-1}+a'_1x^{n-2}+\cdots+a'_{n-1},
\]
then multiplication by $x-\alpha_n$ gives
\begin{equation}
\begin{aligned}
 a_1&=a'_1-\alpha_n,\\
 a_2&=a'_2-\alpha_n a'_1,\\
 a_3&=a'_3-\alpha_n a'_2,\\
&\cdots\\
 a_\nu&=a'_\nu-\alpha_n a'_{\nu-1},\\
&\cdots\\
 a_n&=-\alpha_n a'_{n-1},
\end{aligned}
\tag{3}
\end{equation}
and from this the correctness of the formulas (2) is immediately evident.

It is important to determine the number of terms which occur in each of the sums (2). We denote the number of terms occurring in the sum $(-1)^\nu a_\nu$ by $B_\nu^{(n)}$; it is the number of combinations without repetition of $n$ elements to the $\nu$th class, that is, of every $\nu$ distinct ones among the $n$ elements. To determine it, suppose first that the $B_{\nu-1}^{(n)}$ combinations to the $(\nu-1)$st class have been formed. From each of these combinations, by adding one at a time of the missing $n-\nu+1$ elements, one can derive $n-\nu+1$ combinations to the $\nu$th class. In this way however every combination is formed $\nu$ times, namely by addition of each of its elements; hence one obtains the relation
\begin{equation}
 B_\nu^{(n)}=B_{\nu-1}^{(n)}{n-\nu+1\over \nu}
\tag{4}
\end{equation}
whereas $B_1^{(n)}$ plainly has the value $n$. Therefore, if one multiplies the equations following from (4),
\begin{equation}
\begin{aligned}
 B_1^{(n)}&=n,\\
 B_2^{(n)}&=B_1^{(n)}{n-1\over 2},\\
&\cdots\\
 B_\nu^{(n)}&=B_{\nu-1}^{(n)}{n-\nu+1\over \nu}
\end{aligned}
\tag{5}
\end{equation}
it follows that:
%CONT>>
